% MORSE supplementary material draft (2026-10-13). Material moved out of the main Method and Empirical Study.
% Uses the same preamble and notation as morse_cvpr_revised.tex.

\section{Router Details}
\label{supp:router}

Both branches of the router in Eq.~\eqref{eq:router} are two-layer perceptrons; the state branch applies
$\operatorname{symlog}(v)=\operatorname{sign}(v)\log(1+|v|)$ to its input. Their output layers start from small random
weights (standard deviation $0.01$), so that different units begin with different expert sets.

\paragraph{Weights without renormalization.}
TopK keeps the $k$ largest routing probabilities and does not renormalize them. While the router is uncertain, the
selected experts therefore receive small weights and the shared expert dominates; a routed expert gains weight only
as the router becomes confident, so that specialization departs gradually from the shared computation.

\paragraph{Learned routing noise.}
During training, a second router with the same inputs predicts a scale
$\sigma_b(t)=\operatorname{softplus}(\cdot)+0.01$, and $\sigma_b(t)\,\epsilon$ with $\epsilon\sim\mathcal{N}(0,I)$ is
added to $h_b(t)$ before the selection, which lets every unit explore other experts. The routing objective uses the
same perturbed probabilities with the noise held fixed. At inference, the noise is removed.

\section{Joint Groups}
\label{supp:groups}

The descriptor of every joint provides nine joint-local features (Table~\ref{tab:features}): the absolute rotation axis in the torso frame (3),
the number of direct child joints, the absolute nominal position, the torque and velocity limits, the width of the
range, and the absolute center of the range. The two limits and the width of the range enter on a logarithmic scale.
The absolute values make the features identical for the left and right counterparts of a joint. The position of the
joint in the torso frame, the robot-wide gains, the mass, and the dimensions of the robot are excluded. Damping,
armature, stiffness, and friction loss are excluded as well: the data are collected with these drawn per episode from
a range shared by all joints of a robot, so they describe the robot rather than the joint. The features are
standardized over all joints of all robots, and $k$-means with $k=8$ (best of 50 $k$-means++ initializations) is run
once before training. The normalized mutual information between the groups and robot identity is 0.40, so the groups
are far from a partition by robot. Table~\ref{tab:groups} lists the eight groups.

\begin{table}[t]
\centering
\footnotesize
\setlength{\tabcolsep}{3pt}
\renewcommand{\arraystretch}{1.15}
\caption{Criteria used to group the joints: nine joint-local features of the joint descriptor. Absolute values make
the left and right counterparts of a joint identical. All features are standardized over the joints of all robots
before $k$-means.}
\label{tab:features}
\begin{tabular}{p{2.0cm} p{4.2cm} l}
\toprule
Feature (dim.) & What it describes & Encoding \\
\midrule
\raggedright Rotation axis (3) & \raggedright Direction about which the joint turns, in the torso frame; it separates
roll ($x$), pitch ($y$), and yaw ($z$) joints & absolute \\
\raggedright Child joints (1) & \raggedright Number of joints that directly follow along the limb, \eg, one for a
thigh and none for a knee & count \\
\raggedright Nominal position (1) & \raggedright Angle of the joint in the default standing pose, \ie, how bent it is
at rest & absolute, rad \\
\raggedright Torque limit (1) & \raggedright Largest torque the motor of the joint can exert, \ie, its strength & log,
Nm \\
\raggedright Velocity limit (1) & \raggedright Fastest the joint can turn & log, rad/s \\
\raggedright Range width (1) & \raggedright Span of angles between the two limits of the joint & log, rad \\
\raggedright Range center (1) & \raggedright Middle of the range, relative to the zero angle of the joint & absolute,
rad \\
\bottomrule
\end{tabular}
\end{table}

% Group g is index g-1 in cross_embodiment_offline_rl/algorithms/morse/v8_joint_groups_rev_k8.json.
\begin{table*}[t]
\centering
\footnotesize
\setlength{\tabcolsep}{4pt}
\renewcommand{\arraystretch}{1.15}
\caption{The eight joint groups. Counterparts of a joint (left and right, front and hind) are listed once. Torque and
velocity limits are medians over the joints of a group, with the range of the torque limit in parentheses.}
\label{tab:groups}
\begin{tabular}{cccl p{3.3cm} p{6.5cm}}
\toprule
Group & Joints & Robots & Type & Characteristics & Members \\
\midrule
1 & 52 & 13 & Roll & \raggedright roll axis; 29\,Nm (1.8--201), 20\,rad/s &
Hip abduction of A1, Go1, Go2, ANYmal~B, ANYmal~C, Barkour~v0, Barkour~vB, and Badger; hip roll of H1, G1, Talos,
NAO, and Cassie; shoulder roll of H1, G1, Talos, and NAO; ankle roll of G1 \\
2 & 20 & 6 & Thigh & \raggedright pitch axis; one child joint; nominal 0.9\,rad; range 5.1\,rad &
Thigh of A1, Go1, Go2, and Barkour~vB, and of the front legs of Barkour~v0; elbow pitch of G1 \\
3 & 26 & 7 & Knee & \raggedright pitch axis; no child joint; nominal 1.5\,rad; range 1.9\,rad &
Calf of A1, Go1, Go2, Barkour~v0, Barkour~vB, and Badger; foot of Cassie \\
4 & 54 & 8 & Pitch, larger motors & \raggedright pitch axis; 75\,Nm (12--403) &
Hip and knee flexion of ANYmal~B and ANYmal~C; thigh of Badger and of the hind legs of Barkour~v0; spine of Badger;
hip, knee, and ankle pitch of H1, G1, and Talos; hip pitch and knee of Cassie; shoulder pitch of H1 and G1; elbow of
H1 and Talos; back pitch of Talos \\
5 & 28 & 3 & Pitch, small servos & \raggedright pitch axis; 3.0\,Nm (0.3--4.1) &
All joints of Bittle; head, shoulder, hip, knee, and ankle pitch of OP3 and NAO; elbow roll of NAO \\
6 & 15 & 4 & Yaw, fast and strong & \raggedright yaw axis; 88\,Nm (18--201), 23\,rad/s &
Hip yaw of H1, G1, and Cassie; torso yaw of H1, G1, and Talos; shoulder yaw of H1 and G1; elbow roll of G1 \\
7 & 22 & 4 & Yaw, slow & \raggedright yaw axis; 3.7\,Nm (1.5--100), 3.9\,rad/s &
Hip yaw of Talos and OP3; shoulder, elbow, and wrist yaw of Talos; head yaw of OP3 and NAO; hip yaw-pitch and elbow
yaw of NAO; coxa of Hexapod \\
8 & 24 & 4 & Roll, slow servos & \raggedright roll axis; 2.7\,Nm (2.0--8.0), 1.8\,rad/s &
Shoulder, hip, and ankle roll and elbow of OP3; ankle roll of NAO; wrist roll of Talos; femur and tibia of Hexapod \\
\bottomrule
\end{tabular}
\end{table*}

\paragraph{Reading the table.}
Every joint of the 16 robots is revolute: it rotates about a single axis. We express this axis in the frame of the
torso, in which $x$ points forward, $y$ sideways, and $z$ up. A \emph{roll} joint rotates about the forward axis $x$
and tilts a limb sideways, as hip abduction does when it spreads a leg away from the body. A \emph{pitch} joint
rotates about the sideways axis $y$ and swings a limb forward and backward, as the hip, knee, and ankle do in walking.
A \emph{yaw} joint rotates about the vertical axis $z$ and turns a limb, the torso, or the head to the left or right.
The axis separates the joints most clearly: Groups~1 and~8 hold roll joints, Groups~2--5 pitch joints, and Groups~6
and~7 yaw joints. Joints that perform the same motion therefore meet in one group regardless of the robot; Group~1,
for example, joins the hip abduction of all eight quadrupeds with the hip and shoulder roll of the humanoids.

Within one axis, the remaining features divide the joints further. The number of child joints counts the joints that
directly follow along the limb: the thigh joint of a quadruped is followed by the knee, whereas the knee ends the
leg. The nominal position is the angle of the joint in the default standing pose, and the range is the span of
angles between its two limits (1\,rad $\approx 57^\circ$). These features separate the thighs (Group~2: one child
joint, nominal 0.9\,rad, range 5.1\,rad) from the knees (Group~3: no child joint, a strongly bent nominal pose of
1.5\,rad, range 1.9\,rad). The grouping follows the mechanics of a joint rather than its name: the elbow of G1 joins
the thighs because, like a thigh, it is a bent pitch joint followed by one more joint. The torque limit is the
largest torque the motor of a joint can exert and measures its strength; the velocity limit is the fastest the joint
can turn. These limits set the small robots apart, whose joints are driven by servos, compact motors with an
integrated gearbox and controller that exert at most a few newton-metres. Among the remaining pitch joints, Group~4
holds those driven by larger motors (median 75\,Nm, at least 12\,Nm) and Group~5 those of Bittle, OP3, and NAO
(median 3\,Nm, at most 4.1\,Nm). The yaw joints divide likewise into Group~6, with strong and mostly fast motors
(median 88\,Nm and 23\,rad/s), and Group~7, with slow ones (median 3.9\,rad/s), and the roll joints into Group~1 and
Group~8, whose joints are both weak and slow (median 2.7\,Nm and 1.8\,rad/s).

Three assignments deserve a closer look. First, the left and right counterparts of a joint have the same features,
because the features use absolute values, and they always share a group; front and hind legs share a group unless
their joints differ, as for the thighs of Barkour~v0, whose ranges are centered 0.8\,rad apart (Groups~2 and~4).
Second, the legs of the hexapod point sideways from the torso, so its femur and tibia joints, which move the leg up
and down as the pitch joints of a quadruped do, rotate mainly about the forward axis and join the roll joints in
Group~8. Third, small motors alone do not place a joint in Group~8: the hip and shoulder roll of NAO, which are
followed by further joints and have a narrow range like the hip abduction of the quadrupeds, stay in Group~1, whereas
its ankle roll, which ends the leg, falls into Group~8. Overall, most groups span several robots, and the joints of
one robot spread over several groups; G1, for instance, has joints in Groups~1, 2, 4, and~6. Experts aligned with
these groups can thus be shared by similar joints of different robots and differ between dissimilar joints of the
same robot.

\section{MORSE on the URMA Architecture}
\label{supp:urma}

MORSE applies to any policy of the form of Eq.~\eqref{eq:policy}. To compare directly with the results reported on
the benchmark, we instantiate it on URMA~\citep{bohlinger2024onepolicy}, the architecture adopted
by~\citet{abe2026crossembodimentoffline}, and keep every URMA component that does not carry an expert layer. URMA
consists of three stages: (1) a joint encoder and a foot encoder, whose outputs are aggregated with the general
observation by attention over the descriptors; (2) a core that maps the aggregate to the body representation $z$; and
(3) a decoder that predicts the action mean of every joint from $z$, the joint encoding, and an embedding of the joint
descriptor, next to a separate branch for the standard deviation. In Stage (1), we turn the second layer of the joint
encoder into a joint-level expert layer; in Stage (2), every second layer of the core into an embodiment-level expert
layer; and in Stage (3), the decoder into a bank of decoder experts. The aggregation, the foot encoder, the
standard-deviation branch, and the critics remain unchanged (Fig.~\ref{fig:urma}). The main text describes the
policy with joints only for brevity: in URMA, the feet are encoded and aggregated like the joints, and the robot
descriptor $\bar d_r$ also includes the mean foot descriptor and the number of feet.
% [X] Add the layer widths, the expert hidden widths, the router sizes, and the parameter counts here.

\begin{figure}[t]
\centering
\fbox{\rule{0pt}{1.3in}\rule{0.95\linewidth}{0pt}}
\caption{MORSE instantiated on URMA. URMA modules are shown in gray and the expert layers of MORSE in orange. The
aggregation, the foot encoder, the standard-deviation branch, and the critics are unchanged.}
\label{fig:urma}
\end{figure}

\section{Interaction between Joint Updates}
\label{supp:interaction}

Sparse routing changes which updates meet in the same parameters. For a joint-level layer $b$, let $g_u$ be the
gradient that joint $u$ of one transition contributes to the parameters of the layer. In a dense layer, $g_u$ spans
all parameters, so every pair of joints, within and across robots, interacts through all of them. In an expert layer,
$g_u$ is nonzero only on the shared expert and the experts selected by $u$, and the interaction of two joints reduces
to the experts they both use,
\begin{equation}
\bigl\langle g_u,g_v\bigr\rangle=\sum_{e\in\mathcal{S}_u\cap\mathcal{S}_v}G^e_b(u)\,G^e_b(v)\,
\bigl\langle g^e_u,g^e_v\bigr\rangle,
\label{eq:localized}
\end{equation}
where $\mathcal{S}_u$ contains the shared expert and the experts selected by $u$, and $g^e_u$ is the gradient of the
loss of $u$ with respect to expert $e$ at unit weight. Following the measurement
of~\citet{abe2026crossembodimentoffline}, we compare the agreement between the gradients of different robots on the
dense core and decoder of URMA and on the experts of MORSE [X].
